\documentclass[12pt,a4paper]{article}
\usepackage[T1]{fontenc}
\usepackage[utf8]{inputenc}
\usepackage{mathptmx}            % Times text and math
\usepackage{courier}             % Type 1 typewriter font (avoids bitmap fonts)
\usepackage[a4paper,margin=2.5cm]{geometry}
\usepackage{amsmath}
\usepackage{amssymb}
\usepackage{graphicx}
\usepackage{booktabs}
\usepackage{array}
\usepackage{tabularx}
\usepackage{threeparttable}
\usepackage[font=small,labelfont=bf]{caption}
\usepackage[protrusion=false]{microtype}
\usepackage[numbers,sort&compress]{natbib}
\usepackage[hidelinks]{hyperref}
\usepackage{url}
\usepackage[switch,modulo]{lineno}
\renewcommand\linenumberfont{\normalfont\scriptsize\rmfamily}
\setlength{\parskip}{3pt}
\setlength{\emergencystretch}{3em}
\renewcommand{\arraystretch}{1.12}
\newcommand{\dd}{\mathrm{d}}
\newcommand{\e}[1]{\times10^{#1}}
\newcommand{\rhob}{\bar\rho_L}
\DeclareUrlCommand\file{\urlstyle{tt}}
\makeatletter\g@addto@macro\UrlBreaks{\do\_\do-}\makeatother
\newcommand{\filec}[1]{\texttt{\detokenize{#1}}}
\urlstyle{same}

\title{\textbf{The Time Attraction Law (TAL), Release V20:\\
A Conditional Design Prediction for a\\
Coil Field-Energy Test with Twin Quartz Oscillators}}
\author{David Szabolcsi\\[2pt]
\small Independent experimental researcher, Unterterzen, Switzerland\\
\small ORCID: \href{https://orcid.org/0009-0005-6300-074X}{0009-0005-6300-074X}}
\date{}

\begin{document}
\maketitle
\linenumbers

\begin{abstract}
\noindent
The Time Attraction Law (TAL) postulates that the local rate of proper time depends on the
top-hat averaged mass(-energy) density $\rhob$ through $f_L=\exp(-\gamma_L\rhob)$ with
$\gamma_L=G L^2/c^2$, so that two clocks in environments differing by $\Delta\rhob$ acquire a fractional
frequency difference $\Delta f/f\approx\gamma_L\Delta\rhob$ in magnitude. An extended form of the law
includes the mass equivalent $B^2/(2\mu_0c^2)$ of magnetic field energy. This paper converts that
statement into a quantitative, fully documented design prediction for one fixed laboratory
configuration specified by the author: eight vertical 2000-turn copper coils driven at 14\,mA each,
a main quartz oscillator (OCXO\,8663 type) inside a granite shell and a twin reference oscillator 1.60\,m away.
The finite-solenoid field (current-sheet model) is $11.632900\,\mu$T at the main clock and
$0.430539\,\mu$T at the reference. The differential field energy inside spheres of radius $L$ is
obtained with Gauss--Legendre radial panels, shared Fibonacci angular points and a direct
$r=L/t$ substitution for the outer domain ($L\ge2.6$\,m). The largest calculated differential TAL
signal of the isolated, polarity-even coil-$B^2$ term is $\Delta f/f\approx-1.20\times10^{-48}$
($-1.2017\times10^{-48}$ at $L=0.630957$\,m on the archived grid); an independent Sobol volume
integral gives $-1.2008435\times10^{-48}$ ($-0.0722\,\%$). Convergence, continuity and
outer-integral checks, including an independent third-round review, are reported together with
their remaining limitations. The result is a calculated order of magnitude within an idealised model,
\emph{not} a measured instrument sensitivity and \emph{not} a TAL detection: no ON/OFF/$\pm I$ phase
time series exists yet. A supplementary analysis of SPARC rotation curves contained in the
repository (123-galaxy locked blind subsample) shows that the TAL factor at the published scales is
numerically identical to the pure-baryon baseline.
\end{abstract}

\noindent\textbf{Keywords:} time attraction law; density-dependent proper time; top-hat averaged density;
magnetic field energy; finite solenoid; quartz oscillator comparison; numerical quadrature;
SPARC rotation curves

%=====================================================================
\section{Introduction}\label{sec:intro}

In general relativity the rate of an ideal clock is set by the metric, and in the weak-field limit by
the Newtonian potential, $\dd\tau\approx\dd t\,(1+\Phi/c^2)$ \cite{Will2014}. Optical clocks now resolve
this effect over laboratory height differences of tens of centimetres \cite{Chou2010} and even across
a millimetre-scale atomic sample \cite{Bothwell2022,Zheng2022}, and clock comparisons are routinely
used to search for physics beyond the standard model \cite{Safronova2018,Hees2016}. A separate line of
work has investigated scalar fields whose effective behaviour depends on the local matter density,
such as chameleon models \cite{Khoury2004,Burrage2018}.

The Time Attraction Law (TAL) \cite{TAL_formalism,TAL_publikacio} is a phenomenological proposal of
this broad type: it postulates that the local proper-time rate is an explicit function of the
density averaged over a sphere of radius $L$. In its extended form, the effective density includes the
mass equivalent of magnetic field energy \cite{TAL_publikacio}. The law therefore makes, in principle,
a laboratory prediction for any configuration in which the magnetic energy around one clock can be
switched while a second clock serves as reference.

The purpose of this paper is narrow and deliberately conservative. For one fixed apparatus specified by
the author, it (i)~states the TAL relations used, with exact equation and page references to the source
documents; (ii)~computes the differential TAL frequency shift produced by the coil field energy; (iii)~documents
the numerical method and its verification, including an independent third-round review
\cite{TAL_V20}; and (iv)~summarises, without extension, the SPARC rotation-curve comparison that exists in
the project repository \cite{TAL_repo,TAL_V19,Lelli2016}. All numbers are taken from the archived
calculation package (Zenodo V20, \cite{TAL_V20}; GitHub V18, \cite{TAL_repo}); file names are given so that
each value can be traced.

The central result is a \emph{conditional design prediction}: a calculated order of magnitude within an
ideal current-sheet model. It is not a measurement, not an instrument sensitivity and not evidence for
or against TAL. No clock phase data recorded under switched coil states exist at the time of writing.

%=====================================================================
\section{Theoretical framework}\label{sec:theory}

\subsection{Definitions}
The following definitions are quoted from the TAL formalism document \file{Time_Attraction.pdf}
\cite{TAL_formalism} (copy in the repository root \cite{TAL_repo}). The top-hat averaged density at a
point $\mathbf x$ is (Def.~1, Eq.~(1), p.~3)
\begin{equation}
\rhob(\mathbf x)=\frac{3}{4\pi L^3}\int_{|\mathbf x'-\mathbf x|\le L}\rho(\mathbf x')\,\dd^3x' ,
\label{eq:tophat}
\end{equation}
with $\rhob\to\rho$ for $L\to0$. The coupling constant is (Def.~2, Eq.~(2), p.~3)
\begin{equation}
\gamma_L=\frac{G_N L^2}{c^2}\qquad[\gamma_L]=\mathrm{m^3\,kg^{-1}},
\label{eq:gamma}
\end{equation}
and the time factor is (Def.~3, Eq.~(7), p.~4)
\begin{equation}
f_L(\rhob)=\exp(-\gamma_L\rhob),\qquad \dd\tau_L=f_L\,\dd t,\qquad f_L\in(0,1].
\label{eq:fL}
\end{equation}
The source gives the weak-field clock prediction $\Delta f/f=\gamma_L\Delta\rhob$ (Eq.~(13), p.~5) and
an upper bound $\gamma<10^{-22}\,\mathrm{m^3\,kg^{-1}}$, equivalently $L<367$\,m, obtained from
$\Delta f/f<10^{-18}$ with $\Delta\rho=10^4\,\mathrm{kg\,m^{-3}}$ (Eq.~(3), p.~3). The same bound appears as
$\gamma_{\rm lab}<1.0\times10^{-22}\,\mathrm{m^3\,kg^{-1}}$ in \file{Idovonzas_x5f_publikacio.pdf} (Eq.~(19), p.~3)
\cite{TAL_publikacio}. This bound is the author's own estimate from a generic clock-stability
argument; it is not a constraint reported by a dedicated density-modulation experiment.

\subsection{Field-energy term}
The extended model of \file{Idovonzas_x5f_publikacio.pdf} (Sect.~7, Eqs.~(15)--(16), p.~2)
\cite{TAL_publikacio} adds the mass equivalent of energy densities to the effective density,
\begin{equation}
\rho_{\rm tot}=\rho_b+\rho_{\rm DM}+\frac{B^2}{2\mu_0c^2}+\dots,\qquad
f(\mathbf r,t)=\exp[-\gamma_L\,\rho_{\rm tot}(\mathbf r,t)],
\label{eq:rhotot}
\end{equation}
where the omitted term is the curvature contribution given in the source. Only the magnetic term
$\rho_B=B^2/(2\mu_0c^2)$ is varied in the configuration studied here.

\subsection{Differential prediction for two clocks}
For a static coil current $I$, the field energy enclosed by a sphere of radius $L$ centred on a clock at
$\mathbf x$ is
\begin{equation}
W(<L;\mathbf x)=\int_{|\mathbf x'-\mathbf x|<L}\frac{|\mathbf B(\mathbf x')|^2}{2\mu_0}\,\dd^3x' ,
\end{equation}
so that, by Eq.~\eqref{eq:tophat}, $\bar\rho_{B,L}(\mathbf x)=W(<L;\mathbf x)/(c^2V_L)$ with
$V_L=4\pi L^3/3$. With $\Delta W(L)=W(<L;\mathbf x_{\rm main})-W(<L;\mathbf x_{\rm ref})$ the signed
differential prediction used throughout is
\begin{equation}
\left(\frac{\Delta f}{f}\right)_{\rm diff}=\operatorname{expm1}\!\left[-\gamma\,\Delta\bar\rho_{B,L}\right]
=\operatorname{expm1}\!\left[-\gamma\,\frac{\Delta W(L)}{c^2V_L}\right]
\;\approx\;-\frac{3G_N}{4\pi c^4}\,\frac{\Delta W(L)}{L}\quad(\gamma=\gamma_L),
\label{eq:diff}
\end{equation}
where the last form follows algebraically from Eq.~\eqref{eq:gamma}. A negative value means that the
main clock runs slower than the reference. Because $\mathbf B\propto I$, $\Delta W\propto I^2$; the
package reports $\Delta W(L)/I^2$ in $\mathrm{J\,A^{-2}}$. Two readings of the coupling are evaluated:
(a)~$\gamma=\gamma_L=G_NL^2/c^2$, with $L$ scanned over 81 logarithmically spaced values between 0.01
and 100\,m, and (b)~a fixed $\gamma=10^{-22}\,\mathrm{m^3\,kg^{-1}}$ as a cross-check at the clock bound.
For each reading the largest $|\Delta f/f|$ over the $L$ grid is quoted. The grid maximum is not
claimed to be an accurately determined continuous maximum location.

\subsection{Polarity reversal and background field}\label{sec:polarity}
The pure coil-$B^2$ term is even under polarity reversal. In a constant background field, the change
in energy density also contains a $2\mathbf B_{\rm bg}\cdot\mathbf B_{\rm coil}$ cross term; ideally,
averaging equal-magnitude $\pm I$ states relative to OFF cancels this term. The reported
$-1.20\times10^{-48}$ value refers to this isolated even term, not to the full prediction for every
possible simple ON/OFF measurement. Explicitly, for an unchanged background field $\mathbf B_{\rm bg}$
and coil field $\pm\mathbf b$,
\begin{equation}
\tfrac12\left\{\left[|\mathbf B_{\rm bg}+\mathbf b|^2-|\mathbf B_{\rm bg}|^2\right]+
\left[|\mathbf B_{\rm bg}-\mathbf b|^2-|\mathbf B_{\rm bg}|^2\right]\right\}=|\mathbf b|^2 ,
\label{eq:pm}
\end{equation}
within the linear frequency-shift approximation \cite{TAL_V20}. The calculation does not establish an
actual background-field measurement, $\pm I$ symmetry or a complete prediction for a simple ON/OFF
experiment; those require experimental data.

%=====================================================================
\section{Methodology}\label{sec:method}

\subsection{Experimental geometry}\label{sec:geometry}
All geometry and electrical inputs are author-supplied values recorded in \file{report_hu.md}
(Sect.~1, 1.1, 1.2) and hard-coded in \file{code/final_fixed.py} of the V20 package \cite{TAL_V20}. They are
summarised in Table~\ref{tab:inputs}.

\begin{table}[t]
\centering
\caption{Fixed configuration (author-supplied inputs unless marked ``derived'').
Source: \filec{report_hu.md} Sect.~1; \filec{code/final_fixed.py}; \filec{csv/fixed_case_numbers.json}.}
\label{tab:inputs}
\small
\begin{tabularx}{\textwidth}{>{\raggedright\arraybackslash}p{4.3cm}X}
\toprule
Item & Value \\
\midrule
Coils & 8, vertical axes, parallel connection, identical current direction; axis circle radius 0.45\,m
(coil-body inner edge at 0.30\,m, outer edge at 0.60\,m) \\
Single coil & length 400\,mm; 2000 turns; 0.15\,mm bare Cu wire; plastic former; flange (body) outer
diameter 300\,mm; measured winding diameter 291.7\,mm (Appendix~\ref{app:meas}) \\
Resistance, current & $R=1789\,\Omega$ per coil (all eight identical); $I=14$\,mA per coil \\
Supply & 24\,V, 30\,Ah battery; $U$ and $I$ logged continuously for 12\,h (author) \\
Derived electrical & $14\,\mathrm{mA}\times1789\,\Omega=25.05$\,V; $24\,\mathrm V/1789\,\Omega=13.4$\,mA;
0.351\,W per coil, 2.81\,W total; source current 0.112\,A \\
Derived wire data & wire length $N\pi D=1832.8$\,m per coil; $R_{20}$ from geometry 1788.160\,$\Omega$
($-0.0470\,\%$ vs.\ 1789\,$\Omega$); Cu mass (8 coils) 2.303\,kg \\
Main clock & OCXO, 10\,MHz, SC-cut 3rd overtone; centred, 300\,mm above the bottom plane of the coils;
housing 58--62\,$^\circ$C (author) \\
Granite shell & uniform cylindrical shell, $r=0.09$--$0.22$\,m, $h=0.50$\,m, open, centred on the clock;
2700\,kg\,m$^{-3}$; derived mass 170.9\,kg; 8\,cm gap to the coil bodies \\
Faraday capsule & cylindrical; inner $\varnothing$\,10\,cm, outer $\varnothing$\,16\,cm, height 24\,cm;
layers (outside to inside): 10-layer Al foil, 20\,m of 1.5\,mm bare Cu wire, Al foil, Cu coating,
expanded-polystyrene (Hungarocell) sleeve, 1.2\,mm cast Al housing; 1\,cm gap to the granite \\
Reference clock & full twin OCXO in an identical capsule, same height and temperature, no granite;
1\,m from the outer coil-body surface, i.e.\ 1.60\,m from the main clock along one coil axis direction \\
Signal logging & PicoScope 2205A and Rigol DS1202ZE oscilloscopes (author) \\
\bottomrule
\end{tabularx}
\end{table}

The oscillators are specified by the author as OCXO 8663XS (10\,MHz, 12\,V, SC-cut, 3rd overtone).
A Micro Crystal 8663XS datasheet could not be located; the package therefore compares the author's
parameters with the Oscilloquartz OCXO\,8662/8663 technical specification (Edition March 2003)
\cite{OSA8663}: short-term stability $<1\times10^{-11}$ (0.2--10\,s), ageing $2\times10^{-10}$/day,
temperature stability $<4\times10^{-9}$ p-p, operating range $-20$ to $+70\,^\circ$C. The datasheet
gives no magnetic sensitivity. Whether the operating range refers to the housing temperature is not
stated in the datasheet, so the comparison with the measured 58--62\,$^\circ$C is conditional.

\subsection{Finite-solenoid field}\label{sec:field}
Each coil is modelled as an ideal cylindrical current sheet of radius $a=D/2=0.14585$\,m, length
$\ell=0.400$\,m and $NI$ ampere-turns. The off-axis field uses the Derby--Olbert closed form in terms of
Bulirsch's generalised complete elliptic integral \cite{Derby2010} (\file{code/fields.py}, 12 iterations of
the \texttt{cel} recursion). The coil centres lie on a circle of radius 0.45\,m, 0.1\,m below the clock
plane, so that the clock is 300\,mm above the bottom plane of the coils. The field is the vector sum of the
eight coils. As an independent check, \file{V18/make_revision3_checks.py} sums $8\times2000$ discrete circular
loops (complete elliptic integrals $K$, $E$ from SciPy \cite{SciPy2020}); the reviewer additionally compared the
package field at five points with a separate circular-loop integral \cite{TAL_V20}.

\subsection{Top-hat finite-source integration}\label{sec:tophat}
\file{code/tophat_finite_source.py} implements Eqs.~\eqref{eq:tophat}--\eqref{eq:fL} for finite sources
(constants $G=6.67430\times10^{-11}\,\mathrm{m^3\,kg^{-1}\,s^{-2}}$, $c=299\,792\,458$\,m\,s$^{-1}$ as
in \cite{TAL_formalism}), including the ball volume $V_L$ and the exact ratio
$\operatorname{expm1}(-\gamma\Delta\rhob)$. For static masses (copper windings, granite) the enclosed mass
fraction of a sphere of radius $L$ around each clock is estimated with shared Monte Carlo samples of the
source volume (400\,000 samples for the thin copper shells, $10^6$ for the granite), and
$\Delta\rhob=M\,(\phi_{\rm main}-\phi_{\rm ref})/V_L$ is formed for the same 81 $L$ values. The capsules are
identical for both clocks, so their differential static term is zero by construction.

\subsection{Differential field-energy integrator}\label{sec:integrator}
\file{code/energy_integrator.py} computes $\Delta W(L)$ for $I=1$\,A without radial interpolation.
The shell function $S(r)=4\pi r^2\langle|\mathbf B|^2\rangle_{S^2}/(2\mu_0)$ is evaluated at each radius
by averaging over a Fibonacci spherical point set \cite{Gonzalez2010}. The \emph{same} angular points are
used for both centres at the same $r$ (shared angular points): $n_{\rm hi}=24\,000$ points on shells with
$0.25\le r\le2.4$\,m, which intersect the coils, and $n_{\rm lo}=6000$ points elsewhere. Radial integration
uses 8-point Gauss--Legendre panels with breakpoints at 0, 0.25, 0.75, 2.6 and 3.0\,m (12, 50, 92 and
8 panels per interval in the base setting). Three regimes are used:
\begin{enumerate}
\item $L\le0.95$\,m (inner form; the reference sphere contains no coil):
$\Delta W=W_{\rm main}(<L)-W_{\rm ref}(<L)$.
\item $0.95<L<2.6$\,m (route A, outer difference):
$\Delta W=W_{\rm ref}(>L)-W_{\rm main}(>L)$, with panels from $L$ to 3\,m followed by a tail
$\int_3^\infty$ evaluated with $r=3/t$.
\item $L\ge2.6$\,m (route B): direct outer difference with the substitution $r=L/t$, $t\in(0,1]$,
32-point Gauss--Legendre in $t$ ($n_t=32$), with no truncation and no interpolation.
\end{enumerate}
The outer forms use no separately estimated total energy, so the difference of two independently estimated
total energies no longer enters the outer branch (the earlier $0.01268\,\mathrm{J\,A^{-2}}$ offset near
$L\approx3$\,m is thereby removed \cite{TAL_V20}).

\subsection{Convergence and independent checks}\label{sec:conv}
Four variants were run: $n_{\rm hi}=24\,000$ (base), 6000 and 12\,000, and 12\,000 with doubled radial panel
density; $n_{\rm lo}=6000$ and $n_t=32$ are identical in all four. The spread is reported as the sample
standard deviation $s=\sqrt{\sum(x_i-\bar x)^2/(n-1)}$ with $n=4$. Continuity of routes A and B was tested
at $L=2.6$--$3.01$\,m. As an independent cross-check, a shared-sample volumetric Monte Carlo estimate
($1.5\times10^6$ points per $L$; difference of the two ball indicator functions weighted by the field energy
density) was computed at $L=0.316$, 0.631, 1.0 and 2.0\,m (\file{csv/convergence_mc_shared_samples.csv}).
In the third-round review \cite{TAL_V20}, the main peak was additionally recomputed with a Cartesian Sobol
cubature \cite{Sobol1967} ($4\times2^{19}$ points for the main sphere, $4\times2^{12}$ points for the smooth
reference-sphere integral), and the far-field energy difference with separate Gauss--Legendre polar and radial
quadrature and uniformly spaced azimuthal angles.

\subsection{Clock-comparison reference quantities}\label{sec:clockref}
The author observed that the Lissajous figure of the two OCXOs completes one $360^\circ$ cycle in 1\,h\,5\,min
(3900\,s). Provided this is a true $2\pi$ electrical phase cycle, the conversion is
$\Delta f=1/3900\,\mathrm{Hz}=0.256410$\,mHz and $|\Delta f/f|=2.56410\times10^{-11}$ at 10\,MHz. This is a
\emph{mean beat offset, not a noise floor}. For orientation only, the package also compares the signal with the
cited systematic uncertainty of a Sr lattice clock, $8.1\times10^{-19}$ \cite{Aeppli2024}, with the Sr
two-clock stability $4.8\times10^{-17}$ at 1\,s \cite{Oelker2019}, and evaluates the Sr quadratic Zeeman shift
with the coefficient $-23.0$\,MHz\,T$^{-2}$ \cite{Lu2020}. The Sr value is a cited systematic uncertainty, not a
universal noise floor. Finally, ``main-clock-equivalent'' coefficients are defined as the linear and quadratic
magnetic coefficients $k_1=|\Delta f/f|_{\rm offset}/B_{\rm main}$ and $k_2=|\Delta f/f|_{\rm offset}/B_{\rm main}^2$ that
would ideally produce a frequency change equal in magnitude to the mean beat offset; they are not measured
coefficients and not upper bounds.

\subsection{SPARC rotation-curve data}\label{sec:sparc_method}
The repository contains a supplementary analysis (V19, \cite{TAL_V19}; \file{tal_factor_7th_model.md} and
\file{hetedik_modell_kimenet/} in \cite{TAL_repo}) that applies the TAL factor as a seventh model inside the
hash-locked SPARC baseline runner of the repository's \file{v17_package/github/} directory. The rotation curves are
SPARC \file{Rotmod_LTG.zip} (175 disk galaxies) \cite{Lelli2016}, SHA-256
\texttt{0a80cc90\ldots f588}; the runner (\file{sparc_full_runner.py}) uses a locked blind subsample
\file{blind_123.txt} of 123 galaxies with 2596 radii. The error model is $\sigma=\max(0.05\,V_{\rm obs},20\,\mathrm{km\,s^{-1}})$,
$\chi^2=\sum[(V_{\rm pred}-V_{\rm obs})/\sigma]^2$, $\chi^2_{\rm red}=\chi^2/(N-1)$. The primary TAL model is
$V_{\rm pred}^2=f_L(\bar\rho)\,V_{\rm bar}^2$ with $L_0=1.2\times10^{-11}$\,m, where $\bar\rho$ is a disk-plus-bulge
midplane density proxy ($h_z=0.2$\,kpc, $\Upsilon_{\rm disk}=0.5$, $\Upsilon_{\rm bulge}=0.7$ inside $\bar\rho$ only).
How $f_L$ enters $V(R)$ is an implementation choice not specified in the published material; an inverse coupling
$V^2=V_{\rm bar}^2/f_L$ was run as a sensitivity check. Per-galaxy $\chi^2_{\rm red}$ values were compared with
two-sided Wilcoxon signed-rank and exact sign tests.

%=====================================================================
\section{Results}\label{sec:results}

\subsection{Magnetic field}
Table~\ref{tab:field} lists the field values. The package current-sheet field and the independent
discrete-loop sum agree to $-4.3\times10^{-8}$ (main) and $-7.6\times10^{-9}$ (reference) relative
(\file{V18/revision3_checks.json}); in the reviewer's separate circular-loop test the largest relative
discrepancy over five points was $1.04\times10^{-14}$ \cite{TAL_V20}.

\begin{table}[t]
\centering
\caption{Magnetic field of the eight coils at 14\,mA per coil (current-sheet model, $D=291.7$\,mm).
Sources: \filec{csv/fixed_case_numbers.json} (\texttt{Bpa}, \texttt{Bref\_pa}, \texttt{rows}),
\filec{csv/fixed_case_and_feasibility.csv}, \filec{revision3_review_EN.md} \S2.}
\label{tab:field}
\small
\begin{tabular}{lr}
\toprule
Quantity & Value\\
\midrule
$B$ at main clock & $11.632900\,\mu$T \;(0.8309\,mT\,A$^{-1}$, vertical)\\
$B$ at reference clock (1.60\,m) & $0.430539\,\mu$T \;(30.75\,$\mu$T\,A$^{-1}$)\\
$B$ inside a coil & 71.1\,$\mu$T\\
Magnetic pressure in a coil, $B^2/2\mu_0$ & $2.0\times10^{-3}$\,Pa\\
\bottomrule
\end{tabular}
\end{table}

\subsection{Differential TAL signal}
Figure~\ref{fig:signal} shows the signed differential prediction, Eq.~\eqref{eq:diff}, as a function of $L$.
$\Delta W(L)$ is positive at all evaluated $L$ (no negative large-$L$ tail occurs). With
$\gamma=\gamma_L$ the largest magnitude on the grid is
\begin{equation}
\left(\frac{\Delta f}{f}\right)_{\rm diff}=-1.2017\times10^{-48}\quad\text{at }L=0.630957\ \text{m},
\end{equation}
i.e.\ $\approx-1.20\times10^{-48}$, corresponding to $k=-6.131177672606294\times10^{-45}\,\mathrm{A^{-2}}$ and
$\Delta W/I^2=1.961\,\mathrm{J\,A^{-2}}$ at that $L$ (\file{csv/fixed_case_numbers.json}, \texttt{TAL\_diff};
\file{csv/tal_differential_vs_L.csv}). Accumulated over one day this is $-1.04\times10^{-43}$\,s
($-5.19\times10^{-44}$\,s per 12\,h). With the fixed $\gamma=10^{-22}\,\mathrm{m^3\,kg^{-1}}$ cross-check the
maximum is $-4.63\times10^{-43}$ at $L=0.5623$\,m ($-4.00\times10^{-38}$\,s/day). Selected values are listed in
Table~\ref{tab:signal}.

\begin{figure}[t]
\centering
\includegraphics[width=0.92\textwidth]{fig1_signal_vs_L.pdf}
\caption{Magnitude of the signed differential TAL prediction (main minus reference; all values negative) of
the isolated polarity-even coil-$B^2$ term at 14\,mA per coil, versus averaging radius $L$ (81 grid points).
Solid: $\gamma=\gamma_L=GL^2/c^2$; dashed: fixed $\gamma=10^{-22}\,\mathrm{m^3\,kg^{-1}}$. Dotted vertical lines
mark the inner/outer switch (0.95\,m, where a small matching residual remains, Sect.~\ref{sec:L095}) and the
start of the direct $r=L/t$ outer integral (2.6\,m). Regenerated from \filec{csv/tal_differential_vs_L.csv}.}
\label{fig:signal}
\end{figure}

\begin{table}[t]
\centering
\caption{Selected rows of \filec{csv/tal_differential_vs_L.csv} (base run; values as archived, 4 significant
digits). $\Delta\bar\rho_L$ is main minus reference.}
\label{tab:signal}
\small
\begin{tabular}{rrrrr}
\toprule
$L$ (m) & $\Delta W/I^2$ (J\,A$^{-2}$) & $\Delta\bar\rho_L/I^2$ (kg\,m$^{-3}$\,A$^{-2}$) &
$\Delta f/f$, $\gamma_L$, 14\,mA & $\Delta f/f$, $\gamma=10^{-22}$, 14\,mA\\
\midrule
0.1    & $1.184\e{-3}$ & $3.145\e{-18}$ & $-4.577\e{-51}$ & $-6.164\e{-44}$\\
0.3162 & $6.382\e{-2}$ & $5.361\e{-18}$ & $-7.803\e{-50}$ & $-1.051\e{-43}$\\
0.5623 & 1.582 & $2.363\e{-17}$ & $-1.087\e{-48}$ & $-4.631\e{-43}$\\
0.631  & 1.961 & $2.074\e{-17}$ & $-1.202\e{-48}$ & $-4.065\e{-43}$\\
1      & 2.177 & $5.782\e{-18}$ & $-8.416\e{-49}$ & $-1.133\e{-43}$\\
1.995  & 0.4398 & $1.471\e{-19}$ & $-8.522\e{-50}$ & $-2.883\e{-45}$\\
3.162  & $1.620\e{-3}$ & $1.361\e{-22}$ & $-1.981\e{-52}$ & $-2.667\e{-48}$\\
10     & $2.778\e{-6}$ & $7.378\e{-27}$ & $-1.074\e{-55}$ & $-1.446\e{-52}$\\
100    & $2.636\e{-11}$ & $7.003\e{-35}$ & $-1.019\e{-61}$ & $-1.373\e{-60}$\\
\bottomrule
\end{tabular}
\end{table}

\subsection{Independent check of the main peak}
A targeted rerun of the archived integrator at $L=0.630957$\,m reproduces $-1.2017108\times10^{-48}$, with a
relative difference of $2.2\times10^{-16}$ from the archived value (review \cite{TAL_V20}; the regenerated
\file{V18/revision3_checks.json} gives $-1.2017108238308338\times10^{-48}$ against the JSON value
$-1.2017108238308337\times10^{-48}$). The reviewer's independent volumetric Sobol integral at the same $L$ gives
$-1.2008435\times10^{-48}$, a $-0.0722\,\%$ difference in signal magnitude, with a standard error of
$3.04\times10^{-52}$ estimated from Sobol scrambles. This standard error is not the uncertainty of the full
model: it does not include uncertainties in coil geometry, current, the actual winding or the environment.

\subsection{Convergence}
Table~\ref{tab:conv} reproduces \file{csv/convergence_dW.csv}; Fig.~\ref{fig:conv} shows the four variants
together with the shared-sample Monte Carlo values of Table~\ref{tab:mc}. At the peak ($L=0.631$\,m) the
four-run relative spread is $1.05\times10^{-4}$; the largest relative spread among the tabulated points,
$6.55\times10^{-3}$, occurs at $L=2$\,m, inside the region where the reference-centred shells intersect the coils.

\begin{table}[t]
\centering
\begin{threeparttable}
\caption{Convergence of $\Delta W(L)/I^2$ over four integrator variants ($n_{\rm hi}=24\,000$ base, 6000,
12\,000, and 12\,000 with doubled radial panels). $s$ is the sample standard deviation ($n-1$, $n=4$).
Source: \filec{csv/convergence_dW.csv}.}
\label{tab:conv}
\small
\begin{tabular}{rllll}
\toprule
$L$ (m) & Base value (J\,A$^{-2}$) & Range, min--max (J\,A$^{-2}$) & $s$ (J\,A$^{-2}$) & $s/$mean\\
\midrule
0.1    & 0.00118391 & 0.00118391--0.00118391 & 0\tnote{a} & 0\\
0.3162 & 0.0637737 & 0.0637725--0.0637856 & $6.85\e{-6}$ & $1.07\e{-4}$\\
0.631  & 1.96126 & 1.96091--1.96135 & $2.05\e{-4}$ & $1.05\e{-4}$\\
1      & 2.17691 & 2.17526--2.18903 & $6.53\e{-3}$ & $3.00\e{-3}$\\
2      & 0.426287 & 0.426248--0.431987 & $2.80\e{-3}$ & $6.55\e{-3}$\\
2.6    & 0.00713721 & identical in all 4 runs & 0\tnote{b} & 0\tnote{b}\\
3.162  & 0.00162106 & identical in all 4 runs & 0\tnote{b} & 0\tnote{b}\\
10     & $2.77763\e{-6}$ & identical in all 4 runs & 0\tnote{b} & 0\tnote{b}\\
31.62  & $8.37811\e{-9}$ & identical in all 4 runs & 0\tnote{b} & 0\tnote{b}\\
100    & $2.6365\e{-11}$ & identical in all 4 runs & 0\tnote{b} & 0\tnote{b}\\
\bottomrule
\end{tabular}
\begin{tablenotes}\footnotesize
\item[a] All four runs agree at $L=0.1$\,m; the varied $n_{\rm hi}$ acts only on shells with $0.25\le r\le2.4$\,m.
\item[b] The four runs use identical nlo and nt settings in this region. The zero sample standard deviation is
therefore not a far-field error estimate; far-field convergence must be tested by varying nlo and nt separately.
\end{tablenotes}
\end{threeparttable}
\end{table}

\begin{table}[t]
\centering
\caption{Shared-sample volumetric Monte Carlo check ($1.5\times10^6$ points per $L$).
Source: \filec{csv/convergence_mc_shared_samples.csv}.}
\label{tab:mc}
\small
\begin{tabular}{rrrrrr}
\toprule
$L$ (m) & MC $\Delta W/I^2$ (J\,A$^{-2}$) & MC $1\sigma$ (J\,A$^{-2}$) & $\sigma/\Delta W$ & GL integrator (J\,A$^{-2}$) &
(GL$-$MC)$/\sigma$\\
\midrule
0.3162 & 0.06646 & 0.001444 & 2.17\,\% & 0.06377 & $-1.86$\\
0.631  & 1.949 & 0.0194 & 0.99\,\% & 1.961 & $+0.61$\\
1      & 2.154 & 0.02455 & 1.14\,\% & 2.177 & $+0.92$\\
2      & 0.4266 & 0.01722 & 4.04\,\% & 0.4263 & $-0.02$\\
\bottomrule
\end{tabular}
\end{table}

\begin{figure}[t]
\centering
\includegraphics[width=\textwidth]{fig2_convergence.pdf}
\caption{Left: $\Delta W(L)/I^2$ from the base run (line), the four integrator variants at the check radii
(open symbols) and the shared-sample Monte Carlo estimates with $1\sigma$ error bars (crosses). Right: relative
sample standard deviation over the four variants. Open grey circles are exact zeros obtained with identical
$n_{\rm lo}$, $n_t$ settings and are not far-field error estimates. Regenerated from \filec{csv/convergence_dW.csv},
\filec{csv/convergence_mc_shared_samples.csv} and \filec{csv/tal_differential_vs_L.csv}.}
\label{fig:conv}
\end{figure}

\subsection{\texorpdfstring{Continuity of the outer routes near 3\,m}{Continuity of the outer routes near 3 m}}
Routes A (panels plus $r=3/t$ tail) and B ($r=L/t$) were evaluated on the overlap $L=2.6$--$3.01$\,m
(Table~\ref{tab:cont}). The relative differences at $L=2.99$, 3.00 and 3.01\,m are $0$, $-7.4\times10^{-16}$
and $-1.5\times10^{-15}$; on 2.6--2.9\,m they are at most $2.9\times10^{-15}$ in magnitude. The reviewer's targeted
rerun found the two routes consistent to approximately $10^{-15}$ relative across 2.6--3.01\,m \cite{TAL_V20}.

\begin{table}[t]
\centering
\caption{Route A versus route B for the outer difference. Source: \filec{csv/continuity_test.csv}.}
\label{tab:cont}
\small
\begin{tabular}{rllr}
\toprule
$L$ (m) & Route A (J\,A$^{-2}$) & Route B (J\,A$^{-2}$) & $(A-B)/B$\\
\midrule
2.6  & $7.137211082279\e{-3}$ & $7.137211082279\e{-3}$ & $-2.917\e{-15}$\\
2.7  & $5.182239219640\e{-3}$ & $5.182239219640\e{-3}$ & $-2.008\e{-15}$\\
2.8  & $3.881360582648\e{-3}$ & $3.881360582648\e{-3}$ & $-1.006\e{-15}$\\
2.9  & $2.978621592572\e{-3}$ & $2.978621592572\e{-3}$ & $-4.368\e{-16}$\\
2.99 & $2.387273260006\e{-3}$ & $2.387273260006\e{-3}$ & $0$\\
3.00 & $2.331314942858\e{-3}$ & $2.331314942858\e{-3}$ & $-7.441\e{-16}$\\
3.01 & $2.277039189332\e{-3}$ & $2.277039189332\e{-3}$ & $-1.524\e{-15}$\\
\bottomrule
\end{tabular}
\end{table}

\subsection{Independent outer-integral comparison}
Table~\ref{tab:outer} reproduces the reviewer's comparison of the archived integrator with a separate integral
that uses Gauss--Legendre polar and radial quadrature with uniformly spaced azimuthal angles (not the archived
Fibonacci average) and maps the infinite domain with $r=L/t$ \cite{TAL_V20}. The field function is shared with the
archive. The coarse/fine grids of the separate integral differ by at most a few $10^{-6}$ relative at the nearby points;
this is a convergence test, not a proof of an absolute error bound. At $L=100$\,m the residual decreases with
angular refinement: 6000/12\,000/24\,000 far-field angular points give 2.6364994, 2.6357023 and
$2.6355308\times10^{-11}\,\mathrm{J\,A^{-2}}$, with 32 and then 64 radial $t$ points. No certified error bound is
claimed for every $L$ or for arbitrary distances.

\begin{table}[t]
\centering
\caption{Outer-integral comparison from the independent third-round review
(\filec{revision3_review_EN.md} \S3, \cite{TAL_V20}).}
\label{tab:outer}
\small
\begin{tabular}{rllr}
\toprule
$L$ (m) & Archived integrator $\Delta W$ (J\,A$^{-2}$) & Separate integral (J\,A$^{-2}$) & Rel.\ discrepancy\\
\midrule
2.6      & 0.00713721108 & 0.00713720352 & $+0.000106\,\%$\\
2.818383 & 0.00369084579 & 0.00369083957 & $+0.000169\,\%$\\
3.00     & 0.00233131494 & 0.00233130964 & $+0.000227\,\%$\\
3.162    & 0.00162105726 & 0.00162105265 & $+0.000284\,\%$\\
10       & $2.77763231\e{-6}$ & $2.77754429\e{-6}$ & $+0.00317\,\%$\\
31.62    & $8.37811251\e{-9}$ & $8.37717664\e{-9}$ & $+0.0112\,\%$\\
100      & $2.63649938\e{-11}$ & $2.63556282\e{-11}$ & $+0.0355\,\%$\\
\bottomrule
\end{tabular}
\end{table}

\subsection{\texorpdfstring{Matching residual at $L=0.95$\,m}{Matching residual at L = 0.95 m}}\label{sec:L095}
The inner/outer switch at $L=0.95$\,m is not exactly continuous. The archived inner and outer values are
2.17738947 and 2.17597607\,J\,A$^{-2}$; the matching residual is 0.00141340\,J\,A$^{-2}$, i.e.\ $0.0649\,\%$ of the
inner value. It equals, to $8.9\times10^{-16}$\,J\,A$^{-2}$, the difference between the two total-energy estimates,
2.237861 (main) and 2.236448\,J\,A$^{-2}$ (reference) (\file{csv/fixed_case_numbers.json}, \texttt{jump\_LJ},
\texttt{W\_total}; \file{V18/revision3_checks.json}, \texttt{L095\_residual}), and is attributed to the angular
quadrature error on reference-centred shells that intersect the coils; it affects only the 0.95--2.4\,m range. This
remains a numerical limitation: the curve is not continuous to machine precision, and machine-precision agreement
applies only to the 2.6--3.01\,m A/B overlap.

\subsection{Comparison scales}
Table~\ref{tab:scales} relates the predicted signal to the reference quantities of Sect.~\ref{sec:clockref}. All
ratios are arithmetic comparisons; none is a detection threshold of the present apparatus.

\begin{table}[t]
\centering
\caption{Comparison quantities for the fixed case (14\,mA per coil).
Source: \filec{csv/fixed_case_and_feasibility.csv}, \filec{csv/fixed_case_numbers.json}, \filec{report_hu.md} \S2.}
\label{tab:scales}
\small
\begin{tabularx}{\textwidth}{>{\raggedright\arraybackslash}Xr}
\toprule
Quantity & Value\\
\midrule
Mean beat offset (65-min Lissajous cycle; \emph{mean beat offset, not a noise floor}) & $2.56410\times10^{-11}$\\
$|\Delta f/f|$ / mean beat offset: $\gamma_L$; $\gamma=10^{-22}$ & $4.69\times10^{-38}$; $1.81\times10^{-32}$\\
$|\Delta f/f|$ / datasheet short-term stability $1\times10^{-11}$ \cite{OSA8663} & $1.2\times10^{-37}$\\
$|\Delta f/f|$ / cited Sr systematic uncertainty $8.1\times10^{-19}$ \cite{Aeppli2024} & $1.48\times10^{-30}$\\
Sr white-FM averaging time, $k=2$, $\sigma(1\,\mathrm s)=4.8\times10^{-17}$ \cite{Oelker2019} (algebraic extrapolation only) & $1.6\times10^{57}$\,yr\\
Sr quadratic Zeeman shift at the main-clock field \cite{Lu2020} & $-7.25\times10^{-18}$\\
Field below which the Sr Zeeman shift stays under $8.1\times10^{-19}$ & $<3.9\,\mu$T\\
Main-clock-equivalent coefficient $k_1$ (linear) & $2.20\times10^{-6}\,\mathrm{T^{-1}}$\\
Main-clock-equivalent coefficient $k_2$ (quadratic) & $0.189\,\mathrm{T^{-2}}$\\
Coil current per coil giving $|\Delta f/f|$ = mean beat offset: $\gamma_L$; $\gamma=10^{-22}$ & $6.5\times10^{16}$\,A; $1.0\times10^{14}$\,A\\
\bottomrule
\end{tabularx}
\end{table}

The Sr quadratic Zeeman shift at $11.63\,\mu$T, $-7.25\times10^{-18}$, is quoted as a cited systematic for an
atomic clock in the same field; it is not a universal noise floor and does not apply to quartz, which has no atomic
Zeeman shift. The quoted white-FM averaging time is an algebraic extrapolation of the chosen formula, not an
experimental time prediction. The main-clock-equivalent coefficients only indicate what ideal magnetic frequency
change would equal the observed mean beat offset; without a measured ON/OFF uncertainty and separation of
environmental couplings no upper bound follows from them.

\subsection{Static mass terms}
The granite shell and the copper windings produce static, not modulated, offsets
(\file{csv/fixed_case_numbers.json}, \texttt{granite}, \texttt{cu\_static};
\file{csv/granite_static_differential.csv}, \file{csv/copper_coils_static_differential.csv}). For the granite
(170.9\,kg), the general-relativistic potential term is $\Phi/c^2=-6.28\times10^{-25}$ at the main clock and
$-7.92\times10^{-26}$ at the reference (difference $-5.49\times10^{-25}$); the TAL extreme is
$-9.6\times10^{-26}$ ($\gamma_L$, $L=0.28$\,m) and $-2.1\times10^{-19}$ ($\gamma=10^{-22}$, $L=0.22$\,m). For the
copper (2.303\,kg), the GR terms are $-3.68\times10^{-27}$ and $-1.09\times10^{-27}$ (difference
$-2.60\times10^{-27}$) and the TAL extremes $-6.2\times10^{-28}$ ($\gamma_L$) and $-2.3\times10^{-22}$
($\gamma=10^{-22}$). The capsules are identical, so their differential static term is zero. The mass and material of
the plastic coil formers are unknown; their static differential term could not be calculated.

\subsection{Higher-current cases}
Table~\ref{tab:highI} reproduces the author's list of higher currents. None is feasible with the available supply,
and all exceed the continuous-current limit (0.405\,A) while remaining below the free-air Preece fusing current
(4.65\,A). The Onderdonk estimate is only an idealized adiabatic fusing time; no safe pulse duration follows from
it, and the calculated high-current cases are not operating recommendations. Even in these cases the signal is
$\le5.5\times10^{-44}$ in magnitude.

\begin{table}[t]
\centering
\caption{Higher-current cases (per coil). Source: \filec{csv/fixed_case_and_feasibility.csv}.
Not operating recommendations.}
\label{tab:highI}
\small
\begin{tabular}{rrrrrrr}
\toprule
$I$ (A) & $U=IR$ (V) & $P_8$ (kW) & Onderdonk $t$ (s) & $I/I_{\rm cont}$ & $B_{\rm main}$ (mT) & $\Delta f/f$ ($\gamma_L$)\\
\midrule
1.5 & 2684 & 32.2 & 11.7 & 3.7 & 1.25 & $-1.38\e{-44}$\\
2   & 3578 & 57.2 & 6.6  & 4.9 & 1.66 & $-2.45\e{-44}$\\
3   & 5367 & 128.8 & 2.9 & 7.4 & 2.49 & $-5.52\e{-44}$\\
\bottomrule
\end{tabular}
\end{table}

\subsection{SPARC rotation curves (repository analysis)}\label{sec:sparc}
The repository does \emph{not} contain rotation-curve fit results for all 175 SPARC galaxies. The script
\texttt{Repoduce SPARK 175 v16 1} in the repository root generates synthetic exponential profiles and mock
``observed'' velocities with random scatter when the SPARC directory is absent; it is not an analysis of SPARC data, and
no results from it are used here. The only SPARC analysis with archived results is the seventh-model comparison on the
123-galaxy locked blind subsample (2596 radii) summarised in Table~\ref{tab:sparc}
(\file{hetedik_modell_kimenet/comparison.json}, \file{hetedik_modell_kimenet/run_meta.json},
\file{tal_factor_7th_model.md}; \cite{TAL_repo,TAL_V19}). The six baseline outputs were reported as byte-identical to
\file{v17_package/github/run_output/}.

\begin{table}[t]
\centering
\caption{Median and mean $\chi^2_{\rm red}$ over 123 galaxies (2596 radii) for the six baryonic baselines and the
TAL factor (seventh model). Sources: \filec{hetedik_modell_kimenet/comparison.json} (\texttt{stats},
\texttt{variant\_stats}), \filec{hetedik_modell_kimenet/run_meta.json}, \filec{tal_factor_7th_model.md} \S4.}
\label{tab:sparc}
\small
\begin{tabular}{llrr}
\toprule
Model & $V_{\rm pred}$ & Median & Mean\\
\midrule
vbar & $V_{\rm bar}$ & 3.1433 & 4.2161\\
vbar\_095 & $0.95\,V_{\rm bar}$ & 3.1995 & 4.3725\\
vbar\_105 & $1.05\,V_{\rm bar}$ & 2.9555 & 4.3224\\
no\_bulge & $(V_{\rm gas}^2+V_{\rm disk}^2)^{1/2}$ & 3.3131 & 6.3928\\
no\_gas & $(V_{\rm disk}^2+V_{\rm bul}^2)^{1/2}$ & 4.0474 & 5.2533\\
disk\_only & $|V_{\rm disk}|$ & 4.2623 & 7.5707\\
TAL, $L_0=1.2\times10^{-11}$\,m (primary) & $f_L^{1/2}V_{\rm bar}$ & 3.1433 & 4.2161\\
TAL, $L=367$\,m / $\gamma=10^{-22}$ & forward & 3.1433 & 4.2161\\
TAL, $L=2.36$\,kpc & forward & 3.1433 & 4.2161\\
TAL, fitted $L=2.371$\,Mpc (in-sample, not blind) & forward & 2.7083 & 4.1761\\
TAL, fitted $L=0.473$\,Mpc (in-sample, not blind) & inverse & 3.1425 & 4.3265\\
\bottomrule
\end{tabular}
\end{table}

For the primary model, $\max\gamma_0\bar\rho=6.3\times10^{-67}$ and $\max|f_L-1|=0$ in double precision at all
2596 points; all predicted velocities and all 123 per-galaxy $\chi^2_{\rm red}$ values are identical to the
\texttt{vbar} baseline (123 ties). It therefore ranks exactly where \texttt{vbar} ranks: better than vbar\_095
(97 vs.\ 26 galaxies; Wilcoxon $p=2.57\times10^{-4}$), no\_gas (110 vs.\ 13) and disk\_only (111 vs.\ 12), and
worse than vbar\_105 (32 vs.\ 91; Wilcoxon $p=0.0146$). At $L=2.36$\,kpc, $|f_L-1|\le2.33\times10^{-5}$. The
fitted forward scale lowers the median but is worse than \texttt{vbar} in 81 of 123 galaxies (sign test
$p=5.6\times10^{-4}$) and was obtained in-sample. Order-one scales (tens of Mpc, or $\gamma=10^{23}$) worsen the
per-galaxy fits (\file{tal_factor_7th_model.md} \S4.3). These results are consistent with the source's own statement
that the galactic correction at laboratory-scale $\gamma$ is of order $10^{-34}$ and negligible
(\file{Time_Attraction.pdf}, Sect.~8, Eq.~(14), p.~5 \cite{TAL_formalism}).

%=====================================================================
\section{Discussion}\label{sec:discussion}

\paragraph{Status of the prediction.}
Subject to the numerical limitations and control qualifications listed below, the calculated pure even field-energy
signal of $\approx-1.20\times10^{-48}$ is numerically supported; it is a usable conditional design prediction, not a
validated measurement result \cite{TAL_V20}. It is an acceptable calculated order of magnitude within the chosen
ideal current-sheet model, not a TAL detection and not a measured instrument sensitivity. On the comparison scales
of Table~\ref{tab:scales}, the signal could not be measured with the OCXO pair, nor with a Sr reference clock.

\paragraph{Limitations.}
\begin{enumerate}
\item \emph{No measurement data.} No ON/OFF/$\pm I$ phase or frequency time series, state log, or
switching-correlated thermal measurement exists yet; the dimensional measurement record (Appendix~\ref{app:meas})
does not replace these. The actual sensitivity of the apparatus is therefore undetermined. The next required evidence
is a timestamped phase plus state-label ON/OFF estimate with its uncertainty, not another calculated signal-to-beat
ratio.
\item \emph{Mean beat offset.} The 65-min Lissajous cycle gives a mean beat offset of $2.56410\times10^{-11}$ only if it
is a true $2\pi$ electrical phase cycle. It is a mean beat offset, not a noise floor. Tying the switching schedule to
the beat period is not intrinsically invalid, but the offset does not make 65 minutes optimal; the interval must be
justified through noise and thermal-settling tests.
\item \emph{Far-field convergence.} The four runs use identical nlo and nt settings in this region. The zero sample
standard deviation is therefore not a far-field error estimate; far-field convergence must be tested by varying nlo
and nt separately. A joint refinement ($n_{\rm lo}$ 6000$\to$12\,000, $n_t$ 32$\to$64) changes $\Delta W$ by
$-9.1\times10^{-7}$ to $-3.0\times10^{-4}$ relative at $L=2.6$--100\,m (Appendix~\ref{app:checks}); separate
variation has not yet been performed.
\item \emph{Matching residual at 0.95\,m.} The residual of 0.00141340\,J\,A$^{-2}$ ($0.0649\,\%$ of the inner value)
remains a numerical limitation; the curve is not continuous to machine precision.
\item \emph{Background cross term.} The pure coil-$B^2$ term is even under polarity reversal. In a constant background
field, the change in energy density also contains a $2B_{\rm bg}\cdot B_{\rm coil}$ cross term; ideally, averaging
equal-magnitude $\pm I$ states relative to OFF cancels this term. The reported $-1.20\times10^{-48}$ value refers to this
isolated even term, not to the full prediction for every possible simple ON/OFF measurement.
\item \emph{Thermal control.} Joule heating is also even in the current, so polarity reversal does not separate it. The
proposed control dissipates the same heat load (2.81\,W) in a non-field-generating (bifilar) resistor near the coils
during the OFF intervals. The purpose of the equal-total-power control with a small residual magnetic field is to
reduce switching-correlated thermal changes. Temperature matching between the ON and OFF states and the residual field
of the control must be verified by measurement; any remaining mismatch must be included in the uncertainty of the
frequency estimate. Equal total power does not guarantee identical spatial heat deposition or identical conductive
and radiative coupling to the clocks.
\item \emph{Model idealisation.} The field is that of ideal current sheets with the measured winding diameter; the
actual multilayer winding, current tolerances, magnetic materials, and the environment are not modelled, and none of
the quoted numerical standard errors includes these effects. The winding diameter rests on an author-supplied
measurement record that was not independently certified on site.
\item \emph{Safety.} The Onderdonk estimate is only an idealized adiabatic fusing time; no safe pulse duration follows
from it, and the calculated high-current cases of Table~\ref{tab:highI} are not operating recommendations.
\item \emph{Theory inputs.} The bound $\gamma<10^{-22}\,\mathrm{m^3\,kg^{-1}}$ is the source's own estimate
(Eq.~(3) of \cite{TAL_formalism}). How $f_L$ enters galactic rotation curves is not specified in the published
material; the SPARC comparison rests on stated implementation choices and on a 123-galaxy subsample, not on all 175
galaxies.
\end{enumerate}

\paragraph{Measurement sequence implied by the design.}
The design calls for (1)~a 12\,h baseline without coil current, logging the Lissajous phase of the two OCXOs, battery
voltage and current, and housing temperature; (2)~equal-length ON/OFF segments at 14\,mA for at least 12\,h;
(3)~the same cycle with reversed polarity, so that effects linear in $B$ change sign while $B^2$ effects do not;
(4)~evaluation of ON$-$OFF differences and their uncertainties and of polarity sums and differences, with separate
treatment of temperature and supply couplings; and (5)~the matched-heat-load control above, with only post-settling
segments entering the evaluation (\file{report_hu.md} \S4). Within this sequence the TAL prediction for the isolated
even term is $-1.202\times10^{-48}$.

%=====================================================================
\section{Conclusion}\label{sec:conclusion}
For a fixed, fully specified laboratory configuration, the TAL relations $f_L=\exp(-\gamma_L\rhob)$,
$\gamma_L=GL^2/c^2$, with the magnetic term $B^2/(2\mu_0c^2)$, predict a differential fractional frequency shift of the
isolated, polarity-even coil-$B^2$ term of $\approx-1.20\times10^{-48}$ between a main and a reference oscillator at
14\,mA per coil. This value is reproduced by an independent Sobol cubature to $-0.0722\,\%$, and the previous outer-integral
discrepancy near $L\approx3$\,m has been eliminated at the checked points. The remaining numerical limitations (untested
far-field resolution dependence at fixed nlo/nt, and the 0.0649\,\% matching residual at 0.95\,m) are disclosed. The
result is a conditional design prediction in an ideal current-sheet model, not a measured sensitivity and not a
detection; the next required evidence is an ON/OFF estimate with uncertainty from timestamped phase data and state labels.
In the repository's SPARC analysis, the TAL factor at the published scales is numerically identical to the pure-baryon
baseline.

%=====================================================================
\section*{Data availability}
All inputs, outputs and review texts are archived in Zenodo record V20,
``Time Attraction Law V20 --- audited release with independent third-round review'',
\href{https://doi.org/10.5281/zenodo.23540741}{doi:10.5281/zenodo.23540741} \cite{TAL_V20}; the previous version V19
(seventh-model SPARC analysis) is \href{https://doi.org/10.5281/zenodo.23497128}{doi:10.5281/zenodo.23497128} \cite{TAL_V19}.
The GitHub V18 package, identical in code and CSV/JSON files to the V20 archive and additionally containing the two
text corrections of the third-round review, is available at
\url{https://github.com/unbekanntedavid-glitch/Time-Attraction}, commit
\texttt{042a616596f1204e78ef137da10c262761bb922e} (directory \file{V18/}) \cite{TAL_repo}. SHA-256 checksums are provided
in \file{V18/SHA256SUMS_V18.txt}, \file{V18/SHA256SUMS_ZENODO_V20.txt} and \file{V18/final/SHA256SUMS.txt}. SPARC
rotation curves are available from the SPARC database \cite{Lelli2016}. All materials are licensed under CC BY-NC 4.0.
No clock phase time series exist yet.

\section*{Code availability}
The calculation is reproduced by \file{code/final_fixed.py} with the modules \file{code/fields.py},
\file{code/energy_integrator.py}, and \file{code/tophat_finite_source.py}, written in Python 3 with NumPy \cite{NumPy2020} and
SciPy \cite{SciPy2020}, all in \file{V18/final/}. The computational checks of Appendix~\ref{app:checks} are regenerated by
\file{V18/make_revision3_checks.py}. The SPARC seventh-model scripts and outputs are referenced in
\file{tal_factor_7th_model.md} \cite{TAL_repo,TAL_V19}. Figures~\ref{fig:signal} and \ref{fig:conv} were regenerated from
the archived CSV files with a plotting script supplied with this manuscript source.

\section*{Declarations}
\noindent\textbf{Funding.} No funding was received for this work.\\
\textbf{Competing interests.} The author has no competing interests to declare.\\
\textbf{Author contributions.} D.\,Sz.\ is the sole author: theory, apparatus specification, measurements of the
apparatus dimensions, calculations and writing.\\
\textbf{Ethics approval, consent.} Not applicable.\\
\textbf{Acknowledgements.} The author thanks the independent reviewer of the third-round review archived in
\cite{TAL_V20}.

%=====================================================================
\appendix
\section{Measurement record: winding diameter}\label{app:meas}
\begin{table}[h]
\centering
\caption{Winding-diameter measurement record (author-supplied; \filec{report_hu.md} \S1.2).}
\label{tab:meas}
\small
\begin{tabular}{ll}
\toprule
Item & Value\\
\midrule
Date of measurement & 2025-11-13\\
Measured by & D.\ Szabolcsi\\
Instrument (winding diameter) & caliper, resolution 0.02\,mm\\
Signal measurement & PicoScope 2205A and Rigol DS1202ZE oscilloscopes\\
Readings (5 consecutive) & 291.7; 291.7; 291.7; 291.7; 291.7\,mm\\
Mean & 291.7\,mm\\
Sample standard deviation & 0\,mm (all five readings identical)\\
Location of the 300\,mm dimension & outer diameter of the coil-body flange, not the winding surface\\
\bottomrule
\end{tabular}
\end{table}
\noindent The zero sample standard deviation of five identical readings does not imply zero total measurement uncertainty
of the dimension. Consequences for the geometry: the coil-body inner edge stays at 0.30\,m and the axis circle at
$0.30+0.15=0.45$\,m; the reference is 1\,m from the outer flange surface (0.60\,m), i.e.\ 1.60\,m from the main clock and
1.004\,m from the winding surface; the gap between granite and flange is 8\,cm (8.42\,cm to the winding surface); the gaps
between adjacent coil bodies and windings are 4.4\,cm and 5.3\,cm. The resistance computed from the measured diameter,
1788.160\,$\Omega$, differs from the stated 1789\,$\Omega$ by $-0.0470\,\%$ \cite{TAL_V20}.

\section{Summary of the revision-3 computational checks}\label{app:checks}
\file{V18/revision3_checks.json} was \emph{regenerated by the author's pipeline} on 2026-10-21 by running
\file{V18/make_revision3_checks.py} on the \file{V18/final} package. It is not the reviewer's original file; the reviewer's
script \file{check_revision3.py} is not available. Its contents are computational checks only, with no raw clock data
(Table~\ref{tab:checks}).

\begin{table}[h]
\centering
\caption{Contents of \filec{V18/revision3_checks.json} (author-regenerated).}
\label{tab:checks}
\small
\begin{tabularx}{\textwidth}{>{\raggedright\arraybackslash}p{4.6cm}X}
\toprule
Check & Result\\
\midrule
SHA-256 integrity & 18/18 entries match \filec{V18/final/SHA256SUMS.txt}\\
$B$, main clock & discrete loops $1.1632900301\times10^{-5}$\,T vs.\ package $1.1632899803\times10^{-5}$\,T; rel.\ diff.\ $-4.28\times10^{-8}$\\
$B$, reference clock & $4.3053923117\times10^{-7}$\,T vs.\ $4.3053922792\times10^{-7}$\,T; rel.\ diff.\ $-7.56\times10^{-9}$\\
Main-peak rerun, $L=0.6309573$\,m ($n_{\rm hi}=24\,000$) & $\Delta W/I^2=1.9611445$\,J\,A$^{-2}$; $k=-6.131177672606294\times10^{-45}$\,A$^{-2}$ (identical to JSON); signal $-1.2017108\times10^{-48}$; runtime 56.9\,s\\
Outer integral, $n_{\rm lo}$ 6000$\to$12\,000 and $n_t$ 32$\to$64 jointly &
rel.\ change $-9.08\times10^{-7}$ (2.6\,m), $-1.95\times10^{-6}$ (3.0\,m), $-2.71\times10^{-5}$ (10\,m), $-3.02\times10^{-4}$ (100\,m)\\
Route A vs.\ B at 3.0\,m & rel.\ diff.\ $-7.44\times10^{-16}$\\
Residual at $L=0.95$\,m & 0.0014133960\,J\,A$^{-2}$ ($0.06491\,\%$ of inner value); equals $W_{\rm main}-W_{\rm ref}$ total-energy difference to $8.9\times10^{-16}$\\
\bottomrule
\end{tabularx}
\end{table}
\noindent The joint refinement in the outer-integral row changes nlo and nt together; it does not replace the separate
variation of nlo and nt requested by the review.

%=====================================================================
\begin{thebibliography}{99}\small

\bibitem{Will2014}
C.\,M.\ Will, The confrontation between general relativity and experiment,
Living Rev.\ Relativ.\ \textbf{17}, 4 (2014).
\href{https://doi.org/10.12942/lrr-2014-4}{doi:10.12942/lrr-2014-4}

\bibitem{Chou2010}
C.\,W.\ Chou, D.\,B.\ Hume, T.\ Rosenband, D.\,J.\ Wineland, Optical clocks and relativity,
Science \textbf{329}, 1630--1633 (2010).
\href{https://doi.org/10.1126/science.1192720}{doi:10.1126/science.1192720}

\bibitem{Bothwell2022}
T.\ Bothwell, C.\,J.\ Kennedy, A.\ Aeppli, D.\ Kedar, J.\,M.\ Robinson, E.\ Oelker, A.\ Staron, J.\ Ye,
Resolving the gravitational redshift across a millimetre-scale atomic sample,
Nature \textbf{602}, 420--424 (2022).
\href{https://doi.org/10.1038/s41586-021-04349-7}{doi:10.1038/s41586-021-04349-7}

\bibitem{Zheng2022}
X.\ Zheng, J.\ Dolde, V.\ Lochab, B.\,N.\ Merriman, H.\ Li, S.\ Kolkowitz,
Differential clock comparisons with a multiplexed optical lattice clock,
Nature \textbf{602}, 425--430 (2022).
\href{https://doi.org/10.1038/s41586-021-04344-y}{doi:10.1038/s41586-021-04344-y}

\bibitem{Safronova2018}
M.\,S.\ Safronova, D.\ Budker, D.\ DeMille, D.\,F.\ Jackson Kimball, A.\ Derevianko, C.\,W.\ Clark,
Search for new physics with atoms and molecules, Rev.\ Mod.\ Phys.\ \textbf{90}, 025008 (2018).
\href{https://doi.org/10.1103/RevModPhys.90.025008}{doi:10.1103/RevModPhys.90.025008}

\bibitem{Hees2016}
A.\ Hees, J.\ Gu\'ena, M.\ Abgrall, S.\ Bize, P.\ Wolf, Searching for an oscillating massive scalar field as a
dark matter candidate using atomic hyperfine frequency comparisons,
Phys.\ Rev.\ Lett.\ \textbf{117}, 061301 (2016).
\href{https://doi.org/10.1103/PhysRevLett.117.061301}{doi:10.1103/PhysRevLett.117.061301}

\bibitem{Khoury2004}
J.\ Khoury, A.\ Weltman, Chameleon fields: awaiting surprises for tests of gravity in space,
Phys.\ Rev.\ Lett.\ \textbf{93}, 171104 (2004).
\href{https://doi.org/10.1103/PhysRevLett.93.171104}{doi:10.1103/PhysRevLett.93.171104}

\bibitem{Burrage2018}
C.\ Burrage, J.\ Sakstein, Tests of chameleon gravity, Living Rev.\ Relativ.\ \textbf{21}, 1 (2018).
\href{https://doi.org/10.1007/s41114-018-0011-x}{doi:10.1007/s41114-018-0011-x}

\bibitem{TAL_formalism}
D.\ Szabolcsi, Az Id\H{o}vonz\'as T\"orv\'enye: $f_L=\exp(-\gamma_L\bar\rho_L)$ [The law of time attraction:
density-dependent local time scale], Zenodo (2026), file \file{Time_Attraction.pdf}.
\href{https://doi.org/10.5281/zenodo.23066993}{doi:10.5281/zenodo.23066993}

\bibitem{TAL_publikacio}
D.\ Szabolcsi, Az Id\H{o}vonz\'as T\"orv\'enye. A lok\'alis id\H{o}sk\'ala s\H{u}r\H{u}s\'egf\"ugg\'es\'enek
formalizmusa (dated 2025-08-17), file \file{Idovonzas_x5f_publikacio.pdf} in the Zenodo record
\href{https://doi.org/10.5281/zenodo.22984786}{doi:10.5281/zenodo.22984786} (2026).

\bibitem{TAL_V20}
D.\ Szabolcsi, Time Attraction Law V20 --- audited release with independent third-round review, Zenodo, version v20 (2026).
\href{https://doi.org/10.5281/zenodo.23540741}{doi:10.5281/zenodo.23540741}

\bibitem{TAL_V19}
D.\ Szabolcsi, Time-Attraction Law seventh model independent verification, Zenodo (2026).
\href{https://doi.org/10.5281/zenodo.23497128}{doi:10.5281/zenodo.23497128}

\bibitem{TAL_repo}
D.\ Szabolcsi, Time-Attraction, GitHub repository, commit \texttt{042a616596f1204e78ef137da10c262761bb922e} (V18, 2026).
\url{https://github.com/unbekanntedavid-glitch/Time-Attraction}

\bibitem{Lelli2016}
F.\ Lelli, S.\,S.\ McGaugh, J.\,M.\ Schombert, SPARC: mass models for 175 disk galaxies with Spitzer photometry and
accurate rotation curves, Astron.\ J.\ \textbf{152}, 157 (2016).
\href{https://doi.org/10.3847/0004-6256/152/6/157}{doi:10.3847/0004-6256/152/6/157}

\bibitem{OSA8663}
Oscilloquartz S.A., OCXO 8662/8663 Oven Controlled Crystal Oscillator, Technical Specification, Edition March 2003
(Neuch\^atel). Archived copy: \file{V18/final/refs/oscilloquartz_8663_ed03-2003.pdf} in \cite{TAL_repo},
\url{https://github.com/unbekanntedavid-glitch/Time-Attraction/blob/042a616596f1204e78ef137da10c262761bb922e/V18/final/refs/oscilloquartz_8663_ed03-2003.pdf}

\bibitem{Derby2010}
N.\ Derby, S.\ Olbert, Cylindrical magnets and ideal solenoids, Am.\ J.\ Phys.\ \textbf{78}, 229--235 (2010).
\href{https://doi.org/10.1119/1.3256157}{doi:10.1119/1.3256157}

\bibitem{SciPy2020}
P.\ Virtanen et al., SciPy 1.0: fundamental algorithms for scientific computing in Python,
Nat.\ Methods \textbf{17}, 261--272 (2020).
\href{https://doi.org/10.1038/s41592-019-0686-2}{doi:10.1038/s41592-019-0686-2}

\bibitem{Gonzalez2010}
\'A.\ Gonz\'alez, Measurement of areas on a sphere using Fibonacci and latitude--longitude lattices,
Math.\ Geosci.\ \textbf{42}, 49--64 (2010).
\href{https://doi.org/10.1007/s11004-009-9257-x}{doi:10.1007/s11004-009-9257-x}

\bibitem{Sobol1967}
I.\,M.\ Sobol', On the distribution of points in a cube and the approximate evaluation of integrals,
USSR Comput.\ Math.\ Math.\ Phys.\ \textbf{7}, 86--112 (1967).
\href{https://doi.org/10.1016/0041-5553(67)90144-9}{doi:10.1016/0041-5553(67)90144-9}

\bibitem{Aeppli2024}
A.\ Aeppli, K.\ Kim, W.\ Warfield, M.\,S.\ Safronova, J.\ Ye, Clock with $8\times10^{-19}$ systematic uncertainty,
Phys.\ Rev.\ Lett.\ \textbf{133}, 023401 (2024).
\href{https://doi.org/10.1103/PhysRevLett.133.023401}{doi:10.1103/PhysRevLett.133.023401}

\bibitem{Oelker2019}
E.\ Oelker et al., Demonstration of $4.8\times10^{-17}$ stability at 1\,s for two independent optical clocks,
Nat.\ Photonics \textbf{13}, 714--719 (2019).
\href{https://doi.org/10.1038/s41566-019-0493-4}{doi:10.1038/s41566-019-0493-4}

\bibitem{Lu2020}
X.\ Lu, M.\ Yin, T.\ Li, Y.\ Wang, H.\ Chang, An evaluation of the Zeeman shift of the $^{87}$Sr optical lattice clock
at the National Time Service Center, Appl.\ Sci.\ \textbf{10}, 1440 (2020).
\href{https://doi.org/10.3390/app10041440}{doi:10.3390/app10041440}

\bibitem{NumPy2020}
C.\,R.\ Harris et al., Array programming with NumPy, Nature \textbf{585}, 357--362 (2020).
\href{https://doi.org/10.1038/s41586-020-2649-2}{doi:10.1038/s41586-020-2649-2}

\end{thebibliography}
\end{document}
